%======================================================================
% appendix_full.tex
% Maximally-detailed technical-report-depth appendix for the combined
% main-track paper. Meant to be \input at the \appendix point of main.tex.
% All proofs and tables are ported from paper/report/main.tex and
% paper/report/theory_full_checked_v2.tex, and every experimental number
% is traced to a source file (stated under each table). No number is
% invented or changed. See APPENDIX_PORT_NOTES.md for provenance and flags.
%
% Macros this file uses that are NOT in the combined main.tex preamble
% (move these to the preamble; they are defined here so the file compiles
% standalone in the test wrapper):
%   \Vac  \Sep  \Z  \E  \one  \Piinv  \sign  \conv  \dist  \Lip  \Orb
% They are guarded with \providecommand so re-definition is harmless.
%======================================================================

\providecommand{\Z}{\mathbb{Z}}
\providecommand{\E}{\mathbb{E}}
\providecommand{\Vac}{V_{\mathrm{ac}}}
\providecommand{\Sep}{\mathrm{Sep}}
\providecommand{\one}{\mathbf 1}
\providecommand{\Piinv}{\Pi_{\mathrm{inv}}}
\providecommand{\sign}{\mathrm{sign}}
\providecommand{\conv}{\operatorname{conv}}
\providecommand{\dist}{\operatorname{dist}}
\providecommand{\Lip}{\operatorname{Lip}}
\providecommand{\Orb}{\mathcal O}

%======================================================================
\section{Full structural and optimization theory}\label{app:theory}
%======================================================================
This appendix gives, in full, the proofs the main text states or sketches, and
the supporting structural results. It is self-contained: \S\ref{app:ft-setup}
fixes notation, \S\ref{app:ft-margin}--\S\ref{app:ft-lazy} prove the six
load-bearing results in the order announced in the main body. Where the main
text scopes a claim as proposed or conditional (the effective-dimension deficit
of \S\ref{sec:aniso}, the rich-regime KKT conjecture of the optimization
discussion), that scoping is preserved verbatim here.

\subsection{Setup and standing notation}\label{app:ft-setup}
We work with one-dimensional signals $x\in\R^{\dd}$; the two-dimensional case
replaces $\Z_{\dd}$ by $\Z_{d_1}\times\Z_{d_2}$ and the cyclic discrete Fourier
transform (DFT) by its two-dimensional analogue, without change to any
statement. The cyclic group $G=\Z_{\dd}$ acts by the circular shift
$(S_sx)_i=x_{(i-s)\bmod\dd}$; each $S_s$ is an orthogonal permutation. The
\emph{orbit} of $x$ is $\Orb(x)=\{S_sx\}_{s\in G}$, and a class set is
\emph{orbit-closed} if it is closed under all shifts. Let
$\Vinv=\{v:S_sv=v\ \forall s\}$ and $\Vac=\Vinv^{\perp}$. The orthogonal
projector onto $\Vinv$ is $\Piinv=\tfrac1\dd\one\one^\top=ee^\top$ with
$e=\dd^{-1/2}\one$, so $\Vinv=\mathrm{span}(\one)$ is the constant (DC,
zero-frequency) line and $\Vac=\{v:\sum_iv_i=0\}$ is its zero-mean (AC)
complement. Write $\fdc(x)=e^\top x=\dd^{-1/2}\sum_ix_i$ for the unit DC
functional. With the unitary DFT $\hat x_k=\dd^{-1/2}\sum_jx_je^{-2\pi ijk/\dd}$
and $\omega=e^{2\pi i/\dd}$, a shift multiplies coefficient $k$ by the unit phase
$\omega^{sk}$, so $|\hat x_k|^2$ is shift-invariant.

For a score $f$ correct on $(x,y)$, the $\ell_p$ robust radius is
$r_p(f,x)=\inf\{\lVert\delta\rVert_p:\ \sign f(x+\delta)\neq y\}$; for a linear
score $w^\top x+b$ with $yf(x)>0$ this equals $y(w^\top x+b)/\lVert w\rVert_q$
($q$ dual to $p$). The margin at the correct label is
$M(x)=f_y(x)-\max_{j\neq y}f_j(x)$, positive exactly when $x$ is correct;
$\eta=\E[M]$ and $L=\E\lVert\grad M\rVert_q$ are the dataset means on correctly
classified clean points. For finite sets $A,B$ write the hard-margin separation
$\Sep(A,B):=\tfrac12\dist(\conv A,\conv B)$, zero when the hulls intersect. All
statements match the combined main text's notation and environments.

\begin{lemma}[Fixed subspace for cyclic shifts]\label{lem:cyclic-fixed}
For the cyclic shift action on $\R^\dd$,
$\Vinv=\mathrm{span}(\one)$, $\Piinv=\dd^{-1}\one\one^\top=ee^\top$, and
$\Vac=\{v:\one^\top v=0\}$. The shifts are simultaneously diagonalized by the
DFT, $FS_sF^{*}=\mathrm{diag}(\omega^{sk})_{k=0}^{\dd-1}$, each eigenvalue of
modulus one, so $|\widehat{S_sx}_k|^2=|\hat x_k|^2$ for every $k$.
\end{lemma}
\begin{proof}
If $v\in\Vinv$ then $S_1v=v$ reads $v_i=v_{i-1}$ modulo $\dd$, so all
coordinates are equal and $v\in\mathrm{span}(\one)$; conversely every constant
vector is fixed by every $S_s$. The orthogonal projector onto $\mathrm{span}(\one)$
is $\one\one^\top/\lVert\one\rVert_2^2=\dd^{-1}\one\one^\top=ee^\top$, with
complement $\{v:\one^\top v=0\}$. Every $S_s$ is a circulant permutation, hence
diagonalized by the DFT; its eigenvalues have modulus one, so shifting changes
only Fourier phases and preserves each $|\hat x_k|^2$.
\end{proof}

\subsection{Result 1: the invariant linear margin equals projection separation}
\label{app:ft-margin}
We restate Theorem~\ref{thm:A} of the main text identically, prove it in full
(reduction to one-dimensional DC thresholding; max-margin equals half the
inter-interval distance; the $\pm\one/\sqrt\dd$ direction; recovery of the Ge
$1/\sqrt\dd$ collapse), and give the finite-group generalization.

\begin{theorem}[Invariant linear margin equals projection separation; restated]
\label{thm:A}
A linear score $g(x)=w^\top x+b$ is shift-invariant for all $x$ if and only if
$w=c\one$. Consequently, for any finite separable classes $X_+,X_-$, the
invariant $\ell_2$ max-margin equals the separation of the DC-projected data,
\[
\gamma_{\mathrm{inv}}=\tfrac12\,\dist(I_+,I_-)
=\tfrac12\,\dist\big(\conv\Piinv X_+,\ \conv\Piinv X_-\big),
\]
where $I_\pm=[\min_{x\in X_\pm}\fdc(x),\ \max_{x\in X_\pm}\fdc(x)]$, with
maximizing direction $\pm\one/\sqrt\dd$ oriented toward the positive class. This
holds for any finite separable dataset and does not require orbit-closure; when the
DC-projected intervals $I_+,I_-$ overlap the formula gives $\gamma_{\mathrm{inv}}=0$,
i.e.\ full-space separability does not imply the invariant problem is separable.
\end{theorem}
\begin{proof}
\emph{Characterization of invariant normals.} Invariance of $g$ under every
shift means $w^\top S_sx=w^\top x$ for all $x,s$, i.e.\ $(S_{-s}w-w)^\top x=0$
for all $x$, hence $S_{-s}w=w$ for all $s$. The only vectors fixed by every
cyclic shift are the constants (Lemma~\ref{lem:cyclic-fixed}), so
$w\in\Vinv=\mathrm{span}(\one)$, i.e.\ $w=c\one$; conversely every $c\one$ gives
an invariant score.

\emph{Reduction to one-dimensional DC thresholding.} With $w=c\one$ the score
depends on $x$ only through $\fdc(x)=\one^\top x/\sqrt\dd$, because
$w^\top x=c\,\one^\top x=c\sqrt\dd\,\fdc(x)$. Thus the data reduce to their
scalar DC values and the classes to the intervals
$I_\pm=[\min_{X_\pm}\fdc,\max_{X_\pm}\fdc]$ on the line $\mathrm{span}(e)$.

\emph{Max-margin on the line.} The maximum-margin separator of two point sets on
a line places the threshold midway in the gap between their convex hulls (here
the two intervals $I_\pm$); its geometric margin is half the distance between
the intervals, $\gamma_{\mathrm{inv}}=\tfrac12\dist(I_+,I_-)$. Concretely, if
$I_+$ lies to the right of $I_-$ the optimal unit normal is $+e=+\one/\sqrt\dd$
and $\gamma_{\mathrm{inv}}=\tfrac12(\min_{X_+}\fdc-\max_{X_-}\fdc)$; if $I_-$
lies to the right the normal is $-e$; and when the intervals overlap the
separation, hence the invariant margin, is $0$. Identifying the projected line
with $\mathrm{span}(\one)$ isometrically gives the convex-hull form
$\tfrac12\dist(\conv\Piinv X_+,\conv\Piinv X_-)$.
\end{proof}

\begin{corollary}[Recovery of the Ge $1/\sqrt\dd$ collapse]\label{cor:ge-recover}
For the single white/black dot $x_+=e_j$, $x_-=-e_j$ one has
$\fdc(\pm e_j)=\pm1/\sqrt\dd$, so Theorem~\ref{thm:A} gives
$\gamma_{\mathrm{inv}}=1/\sqrt\dd$ (half-gap convention; the full DC gap is
$2/\sqrt\dd$), the collapse of \citet{ge2021shift}. The unconstrained max-margin
of the two unshifted points $\{e_j\},\{-e_j\}$ is $1$, but on the orbit-closed
dataset $\Orb(e_j)$ versus $\Orb(-e_j)$ the unconstrained linear max-margin is
itself $1/\sqrt\dd$, so the $1$-versus-$1/\sqrt\dd$ contrast is a property of the
two-image protocol, not of invariance alone.
\end{corollary}
\begin{proof}
$\fdc(\pm e_j)=e^\top(\pm e_j)=\pm1/\sqrt\dd$ gives the invariant margin
directly. On $\Orb(e_j)$ vs $\Orb(-e_j)$ every positive basis vector must be
classified positive and every negative one negative by a single hyperplane; the
symmetric max-margin solution is $w\propto\one$ with margin $1/\sqrt\dd$, so the
unconstrained and invariant margins coincide there.
\end{proof}

\begin{proposition}[General finite-group invariant margin]\label{prop:finite-group-app}
Let a finite group $G$ act orthogonally on $\R^\dd$ by matrices $\{R_g\}$, with
fixed subspace $V^G=\{v:R_gv=v\ \forall g\}$ and group-average projector
$P_G=\tfrac1{|G|}\sum_{g}R_g$. A linear score $w^\top x+b$ is exactly
$G$-invariant if and only if $w\in V^G$, and the maximum invariant $\ell_2$
margin for finite separable classes is
$\gamma_{\mathrm{inv}}=\tfrac12\dist(\conv P_GX_+,\conv P_GX_-)$.
Theorem~\ref{thm:A} is the cyclic case $P_G=\Piinv$.
\end{proposition}
\begin{proof}
$P_G$ is the orthogonal projector onto $V^G$: group closure gives
$P_G^2=|G|^{-2}\sum_{g,h}R_gR_h=|G|^{-1}\sum_kR_k=P_G$; each $R_g$ orthogonal
gives $R_g^\top=R_{g^{-1}}$ and $P_G^\top=P_G$; its range is exactly $V^G$.
Invariance of $w^\top x$ for all $x$ is equivalent to $R_g^\top w=w$ for all $g$,
hence (as $G$ is a group) to $R_gw=w$, i.e.\ $w\in V^G$. For $w\in V^G$ we have
$w=P_Gw$, so $w^\top x=w^\top P_Gx$: every invariant linear classifier on
$X_\pm$ is an ordinary linear classifier on the projected sets $P_GX_\pm$, with
the same normal and signed margins. The maximum unit-normal margin separating
two finite sets $A,B$ equals $\tfrac12\dist(\conv A,\conv B)$, attained by the
unit vector along the closest-points segment of the hulls: if every point of
$A,B$ is at signed distance $\ge\gamma$ from a unit-normal hyperplane, so is
every convex combination, so the hulls are at distance $\ge2\gamma$; the segment
midpoint realizes it. Apply this to $A=P_GX_+$, $B=P_GX_-$.
\end{proof}

\begin{corollary}[Projection cannot increase the linear max-margin]\label{cor:proj-mono}
With $\gamma_{\mathrm{full}}=\tfrac12\dist(\conv X_+,\conv X_-)$ and
$\gamma_{\mathrm{inv}}$ as above, $\gamma_{\mathrm{inv}}\le\gamma_{\mathrm{full}}$,
because orthogonal projectors are non-expansive:
$\lVert P_Gu-P_Gv\rVert_2\le\lVert u-v\rVert_2$, so taking the infimum over the
hulls preserves the inequality.
\end{corollary}

\subsection{Result 2: the two-sided bracket (a conditional converse)}
\label{app:ft-bracket}
We restate Proposition~\ref{prop:sandwich} of the main text and give the full
proof, with the explicit reaching-path and co-Lipschitz$(\alpha)$ hypotheses,
the exact linear case $\kappa=1$ \citep{hein2017formal}, the $t^{2n+1}$
large-$\kappa$ counterexample, and the note that it completes the missing upper
arm of \citet{tsuzuku2018lipschitz}.

\begin{proposition}[Two-sided bracket; restated]\label{prop:sandwich}
Let $g$ be continuous and $L$-Lipschitz in $\ell_2$ on a convex set $S$
containing $x$, its nearest boundary point, and the path $\gamma$ below, correct
at $x$ with margin $\eta=y\,g(x)>0$, and write $\mathcal B=\{z:g(z)=0\}$.
\emph{(Lower arm)} $r_2(g,x)\ge\eta/L$. \emph{(Upper arm)} if there is a
continuous path $\gamma:[0,T]\to S$ with $\gamma(0)=x$, $g(\gamma(T))=0$, and a
co-Lipschitz (lower-Lipschitz) bound
$|g(\gamma(t))-g(x)|\ge\alpha\lVert\gamma(t)-x\rVert_2$ up to the first boundary
hit, then $r_2(g,x)\le\eta/\alpha$, and necessarily $\alpha\le L$. Hence
\[
\frac{\eta}{L}\ \le\ r_2(g,x)\ \le\ \frac{\eta}{\alpha},\qquad
\kappa:=\frac{L}{\alpha}\ge1,
\]
so $\etaL$ determines the robust radius up to the local bi-Lipschitz condition
number $\kappa$; the certificate is exact ($r_2=\eta/L$) iff $\kappa=1$.
\end{proposition}
\begin{proof}
Convexity of $S$ keeps the relevant segments inside $S$, so the local Lipschitz
and co-Lipschitz constants apply along them.

\emph{Lower arm.} For any $\delta$ with $\lVert\delta\rVert_2<\eta/L$ the
Lipschitz bound gives $|g(x+\delta)-g(x)|\le L\lVert\delta\rVert_2<\eta=|g(x)|$,
so $g(x+\delta)$ keeps the sign of $g(x)$ and no perturbation below $\eta/L$
flips the label; hence $r_2\ge\eta/L$. This is the standard Lipschitz--margin
certificate \citep{hein2017formal,tsuzuku2018lipschitz}.

\emph{Upper arm.} At $t=T$, $g(\gamma(T))=0$ and $|g(x)|=\eta$, so
$\eta=|g(x)-g(\gamma(T))|\ge\alpha\lVert\gamma(T)-x\rVert_2$, giving
$\lVert\gamma(T)-x\rVert_2\le\eta/\alpha$. Provided $\gamma(T)$ is a \emph{transversal} crossing---$g$
changes sign there, rather than merely touching zero tangentially---the
perturbation $\delta=\gamma(T)-x$ is label-flipping and
$r_2\le\lVert\gamma(T)-x\rVert_2\le\eta/\alpha$. Transversality is a genuine
hypothesis: a tangential zero (e.g.\ $g(\gamma(t))=(t-1)^2\eta$) satisfies the
co-Lipschitz bound yet does not flip the label, so it cannot be dropped.

\emph{$\alpha\le L$.} As $g(\gamma(T))=0\neq g(x)$ we have $\gamma(T)\neq x$, and
the same secant obeys
$\alpha\lVert\gamma(T)-x\rVert_2\le|g(\gamma(T))-g(x)|\le L\lVert\gamma(T)-x\rVert_2$,
so $\alpha\le L$ and $\kappa\ge1$. Combining the arms gives the bracket, tight
iff $\kappa=1$. The upper arm is informative only when $\alpha$ is certified
independently of $r_2$: taking $\alpha$ as the secant slope of the minimal path
itself makes $r_2=\eta/\alpha$ a tautology, whereas the linear case and the
power-spectrum surrogate below fix $\alpha$ from the feature structure.
\end{proof}

\begin{corollary}[Linear invariant features: the converse is exact]\label{cor:linear-exact-app}
If $g(x)=w^\top x+b$ with $w\in\Vinv$ (so $g$ reads $x$ only through $\fdc$),
then along the straight path $\gamma(t)=x-t\,w/\lVert w\rVert_2$ one has
$|g(z)-g(x)|=|w^\top(z-x)|\le\lVert w\rVert_2\lVert z-x\rVert_2$ with equality
along $\pm w$, so $L=\alpha=\lVert w\rVert_2$, $\kappa=1$, and the boundary is
reached at distance $\eta/\lVert w\rVert_2$: the exact point-to-hyperplane
radius $r_2=\eta/\lVert w\rVert_2$ of \citet{hein2017formal}.
\end{corollary}

\begin{remark}[The large-$\kappa$ counterexample: $t^{2n+1}$]\label{rem:tpow-app}
There is no \emph{universal} converse. For the shift-invariant scalar feature
$\Psi_n(x)=t^{2n+1}$ with $t=\fdc(x)$, and points $x_+=e$, $x_-=-e$, the signed
feature margin is $\eta=1$. On the segment $\{te:t\in[-1,1]\}$ the Lipschitz
constant is $L_n=\max_{|t|\le1}(2n+1)t^{2n}=2n+1$, so the certificate gives
$\eta/L_n=1/(2n+1)$. But the boundary is $t=0$, reached from $e$ at Euclidean
distance $1$ (perturbations in $\Vac$ do not change $t$), so $r_2=1$ and
$r_2/(\eta/L_n)=2n+1\to\infty$. In the bracket language this is the large-$\kappa$
end: the straight DC path has secant slope $\alpha=1$, the upper Lipschitz is
$L=2n+1$, so $\kappa=2n+1$ and the bracket $[1/(2n+1),\,1]$ correctly contains
$r_2=1$ at its $\eta/\alpha$ end. The failure of a universal converse is exactly
the unboundedness of $\kappa$ over the model class, not the absence of a bracket
for a fixed model.
\end{remark}

\begin{remark}[Completing the Tsuzuku chain]\label{rem:tsuzuku-app}
\citet{tsuzuku2018lipschitz} (their Eq.~5) bound $r_2$ only from below through
successively tighter Lipschitz constants and state the top inequality
(certificate $\le$ true radius) cannot be guaranteed, measuring it empirically as
a $1.8$--$2.4\times$ gap. The upper arm $r_2\le\eta/\alpha$ of
Proposition~\ref{prop:sandwich} is the matching arm that turns their
one-sided certificate into a two-sided bracket. It is deterministic and tied to a
single fixed classifier, its looseness controlled by the explicit $\kappa$
rather than, as in the randomized-smoothing converse of \citet{cohen2019certified},
by a worst case over classifiers. We verify numerically that the true radius lies
strictly inside the bracket with $\kappa\in[1.6,2.4]$ for the power-spectrum
surrogate.
\end{remark}

For completeness we record the certifiable upper arm for the quadratic-invariant
class, where the reaching path, feature co-Lipschitz constant, and head slope are
all known a priori, so the bracket is genuinely two-sided (not tautological).

\begin{lemma}[Radial co-Lipschitz bound for the power spectrum]\label{lem:ps-colip-app}
Let $\Psi(x)=(|\hat x_k|^2)_k$ with the unitary DFT, and for $r_0,\sigma_0>0$ let
$\mathcal C(r_0,\sigma_0,B)=\{x:r_0\le\lVert x\rVert_2\le B,\ \lVert\Psi(x/\lVert x\rVert_2)\rVert_2\ge\sigma_0\}$.
For $x\in\mathcal C(r_0,\sigma_0,B)$ and $t\in[0,1]$,
$\lVert\Psi(x)-\Psi(tx)\rVert_2\ge r_0\sigma_0\lVert x-tx\rVert_2$.
\end{lemma}
\begin{proof}
Two-homogeneity gives $\Psi(tx)=t^2\Psi(x)$, so
$\lVert\Psi(x)-\Psi(tx)\rVert_2=(1-t^2)\lVert\Psi(x)\rVert_2$ while
$\lVert x-tx\rVert_2=(1-t)\lVert x\rVert_2$; the ratio is
$(1+t)\lVert\Psi(x)\rVert_2/\lVert x\rVert_2\ge\lVert x\rVert_2\lVert\Psi(x/\lVert x\rVert_2)\rVert_2\ge r_0\sigma_0$
on the cone, using $\Psi(x)=\lVert x\rVert_2^2\Psi(x/\lVert x\rVert_2)$ and
$1+t\ge1$.
\end{proof}

\begin{proposition}[Certifiable bracket for radial power-spectrum scores]\label{prop:ps-alpha-app}
Let $g=h\circ\Psi$ with $\Psi$ as above, $x\in\mathcal C(r_0,\sigma_0,B)$,
$g(x)>0$, $g(0)<0$; by continuity there is $t_\star\in(0,1)$ with
$g(t_\star x)=0$, set $z=t_\star x$. If
$h(\Psi(x))-h(\Psi(z))\ge h_-\lVert\Psi(x)-\Psi(z)\rVert_2$ for some $h_->0$,
then $r_2(x)\le g(x)/(h_-r_0\sigma_0)$. If in addition $h$ is $h_+$-Lipschitz on
$\Psi(B_2(0,B'))$ with $B'$ large enough that the certified neighborhood of $x$
stays in the ball, then
\[
\frac{g(x)}{2B'h_+}\ \le\ r_2(x)\ \le\ \frac{g(x)}{h_-r_0\sigma_0},
\qquad \kappa\le\frac{2B'h_+}{h_-\,r_0\sigma_0}.
\]
\end{proposition}
\begin{proof}
$\varphi(t)=g(tx)$ is continuous with $\varphi(0)<0<\varphi(1)$, giving
$t_\star$ and the boundary point $z=t_\star x$ on the explicit radial path.
Lemma~\ref{lem:ps-colip-app} and the head secant give
$g(x)=g(x)-g(z)\ge h_-r_0\sigma_0\lVert x-z\rVert_2$, so
$r_2(x)\le\lVert x-z\rVert_2\le g(x)/(h_-r_0\sigma_0)$. For the lower arm,
$|g(u)-g(v)|\le h_+\lVert\Psi(u)-\Psi(v)\rVert_2\le2B'h_+\lVert u-v\rVert_2$ on
the ball (Lemma~\ref{lem:ps-lip-app}), and the standard certificate applies. The
reaching path is the explicit radial path, the feature co-Lipschitz constant is
the explicit cone constant $r_0\sigma_0$, and $h_-$ is a property of the head on
a known feature segment, so all three are available before any attack is run: the
bracket is genuinely two-sided on the cone, not obtained by naming the unknown
secant slope of the true minimal path.
\end{proof}

\subsection{Result 3: scale-invariance of the ratio (gauge-freeness)}
\label{app:ft-scale}
We restate Lemma~\ref{lem:ratiodegen} of the main text, prove it, and record the
corollary that the direct maximizer is the constant classifier.

\begin{lemma}[Scale-invariance of the ratio; restated]\label{lem:ratiodegen}
The per-sample ratio $M(x)/\lVert\grad M(x)\rVert$ is invariant to positive
rescaling of the score, $f\mapsto cf$ for $c>0$, since numerator and denominator
scale together.
\end{lemma}
\begin{proof}
Under $f\mapsto cf$ the margin $M(x)=f_y(x)-\max_{j\neq y}f_j(x)$ scales to
$cM(x)$ and $\grad M(x)$ to $c\grad M(x)$, so
$cM(x)/\lVert c\grad M(x)\rVert=M(x)/\lVert\grad M(x)\rVert$: the ratio, and the
input-space radius it lower-bounds, are unchanged.
\end{proof}

\begin{corollary}[The direct maximizer is a flat, zero-gradient score]\label{cor:ratiodegen-app}
An objective that maximizes $\etaL$ directly (equivalently penalizes
$\lVert\grad M\rVert/M$ over correctly classified points) is blind to the
absolute margin and admits a flat, zero-gradient score ($\grad M\equiv0$) as a
global optimizer, so $\etaL$ is ill-posed as a training objective and well-posed
only as a diagnostic read at a fixed logit scale.
\end{corollary}
\begin{proof}
The penalty $\E[\lVert\grad M\rVert/M]\ge0$ attains its infimum $0$ as
$\grad M\to0$ with the margin held positive---a flat, zero-gradient score (in the
limit a hard step that separates almost everywhere), which is maximally
non-robust. Because the ratio is scale-invariant (Lemma~\ref{lem:ratiodegen}) the
objective cannot distinguish a large-margin robust classifier from such a
degenerate flat one; this is the constant-function degeneracy noted by
\citet{tsuzuku2018lipschitz} in the large-target limit, reached here through
scale-invariance.
\end{proof}

\subsection{Result 4: the nonlinear escape (power-spectrum span and radius)}
\label{app:ft-ps}
The invariant \emph{linear} family is one-dimensional (only $\fdc$); nonlinear
invariants are richer. We prove the power-spectrum span lemma, the conditional
robust-radius corollary, the ball Lipschitz bound, and the exactly-solvable
nonlinear invariant model the main text mentions as the ``nonlinear escape.''

\begin{lemma}[Power-spectrum span]\label{lem:ps-app}
Define the circular cross-correlation $(w\star x)_s=\sum_jw_jx_{j+s\bmod\dd}$ and
$\Phi_w(x)=\tfrac1\dd\sum_s(w\star x)_s^2$. Then for real $w,x$,
$\Phi_w(x)=\sum_k|\hat w_k|^2|\hat x_k|^2$. Varying a \emph{real} filter $w$
spans the nonredundant real power coordinates $|\hat x_0|^2$, $|\hat x_{\dd/2}|^2$
(when $\dd$ is even), and the conjugate-pair sums
$|\hat x_k|^2+|\hat x_{\dd-k}|^2$ for $0<k<\dd/2$; since $x$ is real these exhaust
its power spectrum. A complex filter isolates each individual $|\hat x_k|^2$.
These features are shift-invariant and supported on $\Vac$ for $k\neq0$.
\end{lemma}
\begin{proof}
The unitary-DFT convolution theorem gives
$\widehat{w\star x}_k=\sqrt\dd\,\overline{\hat w_k}\hat x_k$ up to the correlation
sign convention. By Parseval,
$\Phi_w(x)=\tfrac1\dd\sum_s|(w\star x)_s|^2=\tfrac1\dd\sum_k|\widehat{w\star x}_k|^2
=\sum_k|\hat w_k|^2|\hat x_k|^2$. For real $w$, $\hat w_{\dd-k}=\overline{\hat w_k}$,
so $|\hat w_{\dd-k}|^2=|\hat w_k|^2$: a real filter cannot split a conjugate pair
but can isolate it (support $\hat w$ on $\{k,\dd-k\}$ with conjugate values
keeping $w$ real), and can isolate the self-conjugate $k=0$ and, when $\dd$ even,
$k=\dd/2$. Linear combinations give the stated real span; a complex filter with a
single nonzero coefficient isolates $|\hat x_k|^2$.
\end{proof}

\begin{corollary}[Conditional robust radius]\label{cor:C-app}
For an invariant feature map $\Psi$ with achievable head margin
$\eta=\tfrac12\dist(\conv\Psi(X_+),\conv\Psi(X_-))$ and Lipschitz constant $L$,
the induced classifier $g(x)=a^\top\Psi(x)+\beta$ ($\lVert a\rVert_2=1$, chosen as
the hard-margin separator) has $\mathrm{SC}=1$ and $r_2\ge\eta/L$.
\end{corollary}
\begin{proof}
Invariance of $\Psi$ gives $g(R_gx)=g(x)$, so $\mathrm{SC}=1$. For any data point
and admissible $\delta$,
$|g(x+\delta)-g(x)|=|a^\top(\Psi(x+\delta)-\Psi(x))|\le\lVert a\rVert_2\lVert\Psi(x+\delta)-\Psi(x)\rVert_2\le L\lVert\delta\rVert_2$;
if $\lVert\delta\rVert_2<\eta/L$ then $|g(x+\delta)-g(x)|<\eta\le yg(x)$, so the
sign cannot change and $r_2\ge\eta/L$. The content is which $\Psi$ a
shift-invariant architecture realizes (Lemma~\ref{lem:ps-app}); the inequality
itself is the standard Lipschitz--margin bound, a lower-bound certificate with no
general converse (Remark~\ref{rem:tpow-app}).
\end{proof}

\begin{lemma}[Power-spectrum Lipschitz bound on a ball]\label{lem:ps-lip-app}
Let $Q(x)=(|\hat x_k|^2)_k$ with the unitary DFT. On $\lVert x\rVert_2\le B$,
$\lVert Q(x)-Q(y)\rVert_2\le2B\lVert x-y\rVert_2$, and a scalar quadratic
invariant $f_a(x)=\sum_ka_k|\hat x_k|^2$ has
$\lVert\grad f_a(x)\rVert_2\le2B\lVert a\rVert_\infty$.
\end{lemma}
\begin{proof}
For complex $z,z'$, $||z|^2-|z'|^2|\le|z-z'|(|z|+|z'|)$. Applying this
coordinatewise to $Fx,Fy$ and using unitarity of $F$,
$\lVert Q(x)-Q(y)\rVert_2\le\lVert F(x-y)\rVert_2(\lVert Fx\rVert_2+\lVert Fy\rVert_2)
=\lVert x-y\rVert_2(\lVert x\rVert_2+\lVert y\rVert_2)\le2B\lVert x-y\rVert_2$.
For the gradient, $\grad f_a(x)=2F^{*}\mathrm{diag}(a)Fx$ in the real
identification, so $\lVert\grad f_a(x)\rVert_2\le2\lVert a\rVert_\infty\lVert x\rVert_2\le2B\lVert a\rVert_\infty$.
\end{proof}

\begin{proposition}[An exactly solvable nonlinear invariant model]\label{prop:toy-app}
Let $U,V$ be orthogonal subspaces of $\R^\dd$ and let
$X_+=\{u\in U:\lVert u\rVert_2=a\}$, $X_-=\{v\in V:\lVert v\rVert_2=a\}$ (or
finite shift-orbit subsets whose hulls contain $0$). The quadratic invariant
$f(x)=\lVert P_Ux\rVert_2^2-\lVert P_Vx\rVert_2^2$ classifies $X_+,X_-$
correctly, every point has exact $\ell_2$ robust radius $a/\sqrt2$, and no linear
classifier separates the classes.
\end{proposition}
\begin{proof}
For $x\in X_+$, $P_Vx=0$ and $\lVert x\rVert_2=a$, so $f(x)=a^2>0$; the negative
class is symmetric. Decompose a perturbation as $\delta=\delta_U+\delta_V+\delta_W$
with $W=(U\oplus V)^\perp$. The boundary condition $f(x+\delta)\le0$ reads
$\lVert x+\delta_U\rVert_2^2\le\lVert\delta_V\rVert_2^2$; $\delta_W$ cannot help,
so set it to $0$. With $r=\lVert x+\delta_U\rVert_2$, the smallest
$\lVert\delta_U\rVert_2$ is $|a-r|$ (move along the ray through $x$) and the
smallest feasible $\lVert\delta_V\rVert_2$ is $r$, so
$\lVert\delta\rVert_2^2\ge(a-r)^2+r^2$, minimized at $r=a/2$ with value $a^2/2$;
hence the radius is $a/\sqrt2$, attained by $\delta_U=-(a/2)(x/a)$,
$\lVert\delta_V\rVert=a/2$. No affine hyperplane separates the classes because
$0$ lies in both class convex hulls, so the hulls intersect. This is a point with
radius $a/\sqrt2$ under a quadratic invariant that no linear classifier matches:
the ``nonlinear escape.''
\end{proof}

\subsection{Result 5: radial--tangential split and the envelope}
\label{app:ft-split}
These back the main text's claim that $\etaL$ is a diagnostic, not a target
(\S\ref{sec:signflip}). We prove the exact polar identity
(Proposition~\ref{prop:split-app}) and the envelope-via-achieved-margin
proposition, keeping the ``KKT point need not be globally max-margin'' scoping.

\begin{proposition}[Radial--tangential split of the margin gradient; restated]
\label{prop:split-app}
For a bias-free network with degree-one positively homogeneous activations
(ReLU, leaky-ReLU, absolute value, maxout), $M$ is degree-one homogeneous.
Writing $r=\lVert x\rVert_2$, $u=x/r$, $m(u)=M(u)$, the input gradient splits
into orthogonal radial and tangential parts
$\grad M(x)=m(u)\,u+\grad_{\mathbb S}m(u)$ with
$u\perp\grad_{\mathbb S}m(u)$, so
\[
\lVert\grad M(x)\rVert_2=\frac{M(x)}{\lVert x\rVert_2}\sqrt{1+\lVert\grad_{\mathbb S}\log m(u)\rVert_2^2}.
\]
The radial term is Euler's identity $\langle\grad M(x),x\rangle=M(x)$, giving the
floor $\lVert\grad M(x)\rVert_2\ge M(x)/\lVert x\rVert_2$ and the cap
$\etaL\le\sup_i\lVert x_i\rVert_2$. The upper envelope
$\lVert\grad M(x_i)\rVert_2\le\tau\,M(x_i)/\lVert x_i\rVert_2$ holds exactly when
$\lVert\grad_{\mathbb S}\log m(u_i)\rVert_2\le\sqrt{\tau^2-1}$; equivalently
$\tau_i=\sec\angle(\grad M(x_i),x_i)$, distinct from the bracket condition number
$\kappa=L/\alpha$.
\end{proposition}
\begin{proof}
Degree-one homogeneity gives $M(x)=r\,m(u)$. For a perturbation $h=au+z$ with
$z\perp u$, $Dr(x)[h]=\langle u,h\rangle=a$ and $Du(x)[h]=z/r$ (tangent to the
sphere), so
$DM(x)[h]=a\,m(u)+\langle\grad_{\mathbb S}m(u),z\rangle=\langle m(u)u+\grad_{\mathbb S}m(u),h\rangle$
for all $h$; hence $\grad M(x)=m(u)u+\grad_{\mathbb S}m(u)$. Orthogonality gives
$\lVert\grad M(x)\rVert_2^2=m(u)^2+\lVert\grad_{\mathbb S}m(u)\rVert_2^2$, and
$m(u)=M(x)/\lVert x\rVert_2$ with
$\grad_{\mathbb S}\log m=\grad_{\mathbb S}m/m$ gives the displayed form and the
envelope equivalence. The radial component equals the Euler floor because
$\langle\grad M(x),x\rangle=r\langle\grad M(x),u\rangle=r\,m(u)=M(x)$.
\end{proof}

\begin{proposition}[Envelope via the achieved margin]\label{prop:ib-envelope-app}
Let $M_\theta(x_i)$ be the signed margin at $x_1,\dots,x_n\in\R^\dd\setminus\{0\}$,
parameter-homogeneous of degree $p$ ($M_{a\theta}=a^pM_\theta$). Let $C$ satisfy
$C(a\theta)=a^pC(\theta)$ and $\lVert\grad_xM_\theta(x_i)\rVert_2\le B_iC(\theta)$
at differentiability points, and let the limiting gradient-descent direction
$\bar\theta$ make $x\mapsto M_{\bar\theta}(x)$ positively $1$-homogeneous with
$\min_iM_{\bar\theta}(x_i)>0$. With the achieved normalized margin
$\gamma_{\mathrm{IB}}=\min_iM_{\bar\theta}(x_i)/C(\bar\theta)$ and the
normalization $\widehat\theta=(\min_iM_{\bar\theta}(x_i))^{-1/p}\bar\theta$, at
every support point $M_{\widehat\theta}(x_i)=1$ and
\[
\lVert\grad_xM_{\widehat\theta}(x_i)\rVert_2\le\frac{B_i}{\gamma_{\mathrm{IB}}},\qquad
\lVert\grad_{\mathbb S}\log m(u_i)\rVert_2\le\sqrt{\Big(\tfrac{B_i\lVert x_i\rVert_2}{\gamma_{\mathrm{IB}}}\Big)^2-1}.
\]
\end{proposition}
\begin{proof}
$\gamma_{\mathrm{IB}}$ is scale-invariant: $\bar\theta\mapsto a\bar\theta$ scales
$\min_iM_{\bar\theta}(x_i)$ and $C(\bar\theta)$ both by $a^p$. With
$\alpha=(\min_iM_{\bar\theta}(x_i))^{-1/p}$, parameter homogeneity gives
$M_{\widehat\theta}(x_i)=\alpha^pM_{\bar\theta}(x_i)=M_{\bar\theta}(x_i)/\min_jM_{\bar\theta}(x_j)$,
so $\min_iM_{\widehat\theta}(x_i)=1$, and
$C(\widehat\theta)=\alpha^pC(\bar\theta)=1/\gamma_{\mathrm{IB}}$. The
gradient-control hypothesis gives
$\lVert\grad_xM_{\widehat\theta}(x_i)\rVert_2\le B_iC(\widehat\theta)=B_i/\gamma_{\mathrm{IB}}$;
dividing by $m(u_i)=M_{\widehat\theta}(x_i)/\lVert x_i\rVert_2$ in
Proposition~\ref{prop:split-app} gives the spherical form.
\end{proof}

\begin{remark}[Scope of the envelope]\label{rem:envelope-scope}
For a depth-$D$ bias-free network with $1$-Lipschitz $1$-homogeneous activations
and $C(\theta)=(\lVert\theta\rVert_2/\sqrt D)^D$ the hypothesis holds with
$B_i=\sqrt2$, since $\grad_xM=J^\top(e_y-e_{j^\star})$ with
$\lVert e_y-e_{j^\star}\rVert_2=\sqrt2$ and
$\lVert J\rVert_2\le\prod_\ell\lVert W_\ell\rVert_2\le\prod_\ell\lVert W_\ell\rVert_F\le(\lVert\theta\rVert_2^2/D)^{D/2}=C(\theta)$
by arithmetic--geometric-mean. Gradient descent on homogeneous networks
converges to a KKT point of the max-margin program
\citep{lyu2020gradient,ji2020directional}, which for ReLU networks need not be
globally max-margin \citep{vardi2022margin}; the bound therefore uses the
\emph{achieved} normalized margin $\gamma_{\mathrm{IB}}$, recovering the global
form only when the two coincide. The constant cannot be data-independent: for
$(\pm e_1,\pm1),(\pm e_2/A,\pm1)$ in $\R^2$ the hard-margin solution is
$w^\star=(1,A)$, giving $\lVert\grad_{\mathbb S}\log m(e_1)\rVert_2=A\to\infty$;
the linear case is exact with $\tau_i=\lVert x_i\rVert_2/\gamma$
\citep{soudry2018implicit}. Empirically the factor is large, $\tau\approx26$ on
the trained networks, exactly the slack that lets a gradient penalty lower
$\lVert\grad M\rVert$ while the margin falls with it, so $\etaL$ diagnoses
robustness but is not a target.
\end{remark}

\subsection{Result 6: the lazy-regime cluster theorem}\label{app:ft-lazy}
We prove that orbit-averaging is an orthogonal projection onto the invariant
RKHS subspace (non-expansive), so the tied minimum-norm interpolant is never
worse than the unconstrained one; on orbit-closed data the two \emph{coincide},
so any empirical advantage is a realized finite-width effect. We keep the
smooth-activation caveat (bare ReLU NTK not Lipschitz). We also record the
orbit-bundle and orbit-rank lemmas and the degeneracy propositions the main
optimization discussion invokes.

\begin{lemma}[Orbit-bundle equivalence]\label{lem:bundle-app}
A one-hidden-layer circular convolutional network with global average pooling,
$f(x)=\sum_r\alpha_r\tfrac1\dd\sum_{s\in G}\sigma(\langle w_r,S_sx\rangle+b_r)+\beta$,
equals a two-layer fully connected network whose hidden units are tied into
shift-orbit bundles: $f(x)=\sum_r\sum_s\tfrac{\alpha_r}\dd\sigma(\langle S_{-s}w_r,x\rangle+b_r)+\beta$,
each filter contributing the orbit $\{S_{-s}w_r\}$ sharing output weight
$\alpha_r/\dd$ and bias $b_r$. The map is exactly invariant.
\end{lemma}
\begin{proof}
Orthogonality of $S_s$ gives
$\langle w_r,S_sx\rangle=\langle S_s^\top w_r,x\rangle=\langle S_{-s}w_r,x\rangle$,
which substituted into the pooling expression gives the tied two-layer form.
Invariance follows by reindexing the pooling sum: $f(S_tx)$ replaces the sum over
$s$ by a sum over $s'=s+t$ ranging over all of $G$, returning $f(x)$.
\end{proof}

\begin{lemma}[Rank of a shift orbit]\label{lem:rank-app}
For a base pattern $u\in\R^\dd$, the cyclic orbit $\{S_su\}_{s\in\Z_\dd}$ has rank
exactly $\#\{k:\hat u_k\neq0\}$, the number of nonzero Fourier coefficients of
$u$.
\end{lemma}
\begin{proof}
The matrix of orbit columns is circulant, hence diagonalized by the DFT with
eigenvalues $\propto\hat u_k$; multiplication by the invertible DFT preserves
rank, so the rank equals the number of nonzero eigenvalues.
\end{proof}

\begin{proposition}[The canonical orbit bases are degenerate, in opposite ways]
\label{prop:rank-app}
Let each class be the shift orbit of a single base pattern. (i) For a single
cosine of frequency $k$, Lemma~\ref{lem:rank-app} gives orbit rank $2$ (it spans
$\{\cos,\sin\}$ at frequency $k$), violating the near-orthogonal cluster
assumption of \citet{frei2023double,li2025feature}, and the signed orbit mean, the
direction carrying the non-robustness conclusion there, is $\mathbf 0$. (ii) For a
localized spike $e_j$, the orbit is full-rank and orthonormal (the cluster
assumption holds) but the power spectrum is phase-blind, so $\Sep(\Psi)=0$ and the
invariant nonlinear classifier cannot separate the two classes. Neither canonical
base is a valid testbed for the re-pointing hope: in (i) the non-robustness
premise fails, in (ii) the invariant-separation conclusion fails.
\end{proposition}
\begin{proof}
(i) A real cosine of frequency $k\neq0,\dd/2$ has $\hat u$ supported on
$\{k,\dd-k\}$, so by Lemma~\ref{lem:rank-app} the orbit rank is $2$; its columns
span the two-dimensional $\{\cos,\sin\}$ block, far below the $\Omega(\dd)$
near-orthogonal directions the cited non-robustness results need. The signed
orbit mean $\tfrac1\dd\sum_sS_su$ projects onto the DC line, which a zero-mean
cosine annihilates, giving $\mathbf 0$. (ii) $\Orb(e_j)$ is the standard basis
(full rank, orthonormal); but $|\widehat{e_j}_k|^2=1/\dd$ for all $k$, so
$+e_j$ and $-e_j$ (indeed any two spikes) share a power spectrum and no
power-spectrum feature separates them, $\Sep(\Psi)=0$. Numerically for $\dd=64$,
$k=5$ the cosine orbit has rank $2$ and signed mean of norm $1.7\times10^{-16}$,
and the spike orbit has rank $\dd$ and power-spectrum gap $0$.
\end{proof}

\begin{proposition}[The degeneracy is non-generic]\label{prop:generic-app}
If a base $u$ has all Fourier coefficients nonzero, its orbit is full-rank
(Lemma~\ref{lem:rank-app}); and two such bases with distinct power spectra are
power-spectrum separable, with feature margin $\tfrac12\lVert Q(u)-Q(v)\rVert_2$.
Hence generic single-orbit data is simultaneously full-rank and
invariant-separable, so the degenerate corner of Proposition~\ref{prop:rank-app}
does not cover single-orbit data in general.
\end{proposition}
\begin{proof}
Full rank is Lemma~\ref{lem:rank-app}. The power spectrum is shift-invariant
(Lemma~\ref{lem:cyclic-fixed}), so each orbit maps to one point $Q(u)$, $Q(v)$;
if $Q(u)\neq Q(v)$ the maximum-margin separator between two points has margin
half their Euclidean distance. For $u$ drawn from any law absolutely continuous
w.r.t.\ Lebesgue measure, each event $\hat u_k=0$ is a proper linear constraint
of probability zero, and for independent continuous $u,v$ the equalities
$|\hat u_k|^2=|\hat v_k|^2$ for all $k$ hold with probability zero; hence the
generic claim.
\end{proof}

\begin{theorem}[Lazy-regime cluster theorem: no minimum-norm certified-Lipschitz
gap on orbit-closed data]\label{thm:lazy-cluster-app}
Work in the lazy/NTK regime, where each trained network equals the
minimum-RKHS-norm interpolant of its fixed neural tangent kernel
\citep{jacot2018neural}: write $K(x,x')=\langle\Phi(x),\Phi(x')\rangle_{\mathcal H}$
for the tangent kernel with feature map $\Phi$ and RKHS $\mathcal H_K$, and assume
a smooth activation so that $\Phi$ is $L_\Phi$-Lipschitz on the data ball
$\{\lVert x\rVert_2\le B\}$. Let the data be orbit-closed with labels constant on
orbits, let the kernel be shift-invariant ($K(S_sx,S_sx')=K(x,x')$), and let
$\Pi_Gf=\tfrac1\dd\sum_sf(S_s\cdot)$ be the orbit-averaging (Reynolds) operator,
whose range $\mathcal H^G=\Pi_G\mathcal H_K$ is the invariant subspace realized by
the shift-tied orbit bundle of Lemma~\ref{lem:bundle-app}. For the unit-margin
min-norm problem $\min_f\lVert f\rVert_K$ s.t.\ $y_if(x_i)\ge1$, let
$f_{\mathrm{unc}}$ solve over $\mathcal H_K$ and $f_{\mathrm{inv}}$ over
$\mathcal H^G$, and set $L_{\mathrm{unc}}=L_\Phi\lVert f_{\mathrm{unc}}\rVert_K$,
$L_{\mathrm{inv}}=L_\Phi\lVert f_{\mathrm{inv}}\rVert_K$. Then
\begin{enumerate}[leftmargin=1.7em,label=\textup{(\roman*)}]
\item $\Pi_G$ is the orthogonal projection of $\mathcal H_K$ onto $\mathcal H^G$,
hence non-expansive, with
$\lVert f\rVert_K^2=\lVert\Pi_Gf\rVert_K^2+\lVert(I-\Pi_G)f\rVert_K^2$
\citep{elesedy2021provably};
\item $\Pi_Gf_{\mathrm{unc}}$ is feasible for the invariant problem, so
$\lVert f_{\mathrm{inv}}\rVert_K\le\lVert\Pi_Gf_{\mathrm{unc}}\rVert_K\le\lVert f_{\mathrm{unc}}\rVert_K$;
\item therefore $L_{\mathrm{inv}}\le L_{\mathrm{unc}}$, and at the common unit
margin $\eta$ the invariant model has the weakly larger certified ratio,
$\eta/L_{\mathrm{inv}}\ge\eta/L_{\mathrm{unc}}$. Moreover, because the RKHS norm
is strictly convex and the data is orbit-closed, the unconstrained min-norm
interpolant is itself invariant, $f_{\mathrm{unc}}=\Pi_Gf_{\mathrm{unc}}$, so the
orbit-variance $\lVert(I-\Pi_G)f_{\mathrm{unc}}\rVert_K^2$ vanishes and
$L_{\mathrm{inv}}=L_{\mathrm{unc}}$: on orbit-closed data there is \emph{no}
strict minimum-norm gap. The cluster-geometry advantage seen empirically is a
realized-local-Lipschitz / finite-width effect, not a smaller min-norm
certificate.
\end{enumerate}
\end{theorem}
\begin{proof}
(i) For a shift-invariant kernel $\Pi_G$ is self-adjoint and idempotent on
$\mathcal H_K$ with spectrum in $\{0,1\}$, hence the orthogonal projection onto
its fixed-point set $\mathcal H^G$; the displayed identity is the Pythagorean
decomposition for that projection \citep{elesedy2021provably}. (ii) Orbit-closed
data with orbit-constant labels give $y_if_{\mathrm{unc}}(S_sx_i)\ge1$ for all
$s$, so $y_i(\Pi_Gf_{\mathrm{unc}})(x_i)\ge1$: $\Pi_Gf_{\mathrm{unc}}\in\mathcal H^G$
is feasible, whence
$\lVert f_{\mathrm{inv}}\rVert_K\le\lVert\Pi_Gf_{\mathrm{unc}}\rVert_K$, bounded by
$\lVert f_{\mathrm{unc}}\rVert_K$ via (i). (iii) For any $f\in\mathcal H_K$ and
$x,x'$ in the data ball,
$|f(x)-f(x')|=|\langle f,\Phi(x)-\Phi(x')\rangle_K|\le\lVert f\rVert_K\lVert\Phi(x)-\Phi(x')\rVert_K\le L_\Phi\lVert f\rVert_K\lVert x-x'\rVert_2$,
so $f$ is $L_\Phi\lVert f\rVert_K$-Lipschitz there; applying this to
$f_{\mathrm{inv}},f_{\mathrm{unc}}$ and using (ii) gives
$L_{\mathrm{inv}}\le L_{\mathrm{unc}}$, and dividing the common $\eta$ gives the
ratio comparison. The orthogonality identity of (i) makes the norm gap exactly
$\lVert(I-\Pi_G)f_{\mathrm{unc}}\rVert_K^2$. On orbit-closed data each
$S_sf_{\mathrm{unc}}$ is feasible with $\lVert S_sf_{\mathrm{unc}}\rVert_K=\lVert f_{\mathrm{unc}}\rVert_K$,
so were $f_{\mathrm{unc}}$ non-invariant the average
$\Pi_Gf_{\mathrm{unc}}$ would be feasible of strictly smaller norm by strict
convexity, contradicting minimality; hence $f_{\mathrm{unc}}=\Pi_Gf_{\mathrm{unc}}$
and $L_{\mathrm{inv}}=L_{\mathrm{unc}}$.
\end{proof}

\begin{remark}[Smooth-activation caveat; what remains open]\label{rem:lazy-relu-app}
The norm contraction (i)--(ii) is activation-independent, but step (iii) uses
smoothness: the bare two-layer ReLU NTK feature map is \emph{not} Lipschitz (its
RKHS contains unit-norm functions of unbounded Lipschitz constant,
\citealp{bietti2019inductive}), so for ReLU one must take a smooth activation or
carry $\sup_{\lVert x\rVert_2\le B}\lVert\grad\Phi(x)\rVert$ explicitly in place
of $L_\Phi$. Theorem~\ref{thm:lazy-cluster-app} thus proves the cluster
advantage in the lazy, smooth-activation regime; the synthetic stress test's
tied-versus-unconstrained $\etaL$ gap (orbit-variance $\approx0.3$, ratio up to
$4.6\times$ in the lazy regime) is therefore a realized, finite-width effect. The
rich (small-init) regime and the bare-ReLU trajectory, where the fixed-kernel
picture breaks down \citep{chizat2019lazy} and the realized-Lipschitz selection
among equal-norm interpolants is the open question, stay open: this is the
rich-regime KKT conjecture the main text leaves open, on the same wall reached by
\citet{frei2023double,min2024can,li2025feature}.
\end{remark}

\begin{corollary}[Rich two-layer: the max-margin values coincide]\label{cor:no-gap-var-app}
The no-gap conclusion of Theorem~\ref{thm:lazy-cluster-app} needs only that the
norm is $G$-invariant and convex and that Reynolds averaging preserves the
orbit-constant margin constraints. For any $G$-invariant convex norm
$\lVert\cdot\rVert$, Jensen gives $\lVert\Pi_Gf\rVert\le\lVert f\rVert$, and on
orbit-closed data $\Pi_Gf$ is invariant and feasible, so the invariant and
unconstrained unit-margin minimum-norm values coincide. This holds for the
variation (atomic) norm of an infinite-width two-layer network; since the
directional limit of gradient descent on a homogeneous network is a max-margin
(KKT) interpolant \citep{lyu2020gradient,ji2020directional}, which for wide
two-layer networks is the minimum-variation-norm max-margin classifier
\citep{chizat2020implicit}, weight sharing cannot raise that max-margin. The
variation norm is not strictly convex, so an invariant optimum exists but need
not be unique; the open rich-regime question is whether the dynamics select an
invariant, lower-local-gradient optimum among equal-norm interpolants, not
whether a smaller-norm optimum exists.
\end{corollary}

%======================================================================
\section{Full experimental tables and protocols}\label{app:extra}
%======================================================================
Every number below is traced to a source file, named under its table. Where a
number in the vision-paper table could not be reconciled with a JSON it is
flagged in the caption and in \texttt{APPENDIX\_PORT\_NOTES.md} rather than
copied. Correlations use Pearson for the CIFAR/ResNet grids and Spearman for the
per-image and RobustBench panels, matching the main text.

\subsection{Weak-attack table}\label{app:weakattack-full}
The ``invariance helps'' effect is a weak-attack artifact. Table~\ref{tab:weakattack-full}
reports the capacity-matched ResNet arms plus the ported TIPS operator,
\emph{standard}-trained (no adversarial training) and attacked with FGSM, PGD-10,
and AutoAttack at small $\ell_\infty$ budgets, the protocol of the invariance-helps
reports \citep{saha2024improving,wang2025bridging}. Under FGSM at $1/255$ the
exactly invariant (exact-APS) and shift-augmented arms sit above the standard
baseline, reproducing that literature's ordering for the built-in-invariance
arms; adding TIPS's auxiliary losses restores its FGSM advantage too; under
AutoAttack every arm is at or near zero by $\varepsilon=2/255$, so the ordering
ranks models that are all non-robust. FGSM overstates robust accuracy here by more
than a factor of $30$ relative to PGD-10 at the same budget. The rotation-equivariant
architecture of \citet{wang2025bridging}, retrained with their code and recipe,
reproduces their weak-attack numbers (Wang CascadedGCNN: clean $87.2\%$, FGSM
$28.53\%$ and PGD-40 $1.75\%$ at their headline budget $0.03$ in their normalized
convention) yet retains only $0.1\%$ under threat-matched AutoAttack (pixel-space
APGD at $0.015$), and $0.0\%$ at $0.03$.

\begin{table}[t]\centering
\caption{Standard-trained ResNet arms (width $1.0$) under weak and strong attacks;
$\ell_\infty$ budgets in $/255$. AutoAttack on $512$ points (binomial $95\%$
half-width $\approx0.007$ near zero). The FGSM ordering at $\varepsilon{=}1$
reproduces the invariance-helps reading for the built-in-invariance arms
(exact-APS and shift-augmented above standard), while the AutoAttack columns show
every arm is non-robust. Source:
\texttt{paper/results/resnet\_scale/tips\_weakattack\_w1.0.json} (standard,
blurpool, aps, aug, tips) and \texttt{c10stdaux\_tips\_s0.json} (tips+aux);
PGD is $10$-step.}
\label{tab:weakattack-full}
\footnotesize
\begin{tabular}{lcc ccc cc}
\toprule
arm & clean & consist. & FGSM$_1$ & FGSM$_2$ & PGD$_1$ & AA$_2$ & AA$_4$ \\
\midrule
standard          & 0.931 & 0.965 & 0.477 & 0.244 & 0.263 & 0.006 & 0.000 \\
anti-aliased      & 0.935 & 0.978 & 0.472 & 0.262 & 0.220 & 0.002 & 0.000 \\
exact (APS)       & 0.926 & 1.000 & 0.497 & 0.277 & 0.300 & 0.004 & 0.000 \\
shift-aug         & 0.941 & 0.973 & 0.532 & 0.315 & 0.306 & 0.004 & 0.000 \\
TIPS (operator)   & 0.929 & 0.974 & 0.459 & 0.229 & 0.261 & 0.000 & 0.000 \\
TIPS + aux losses & 0.931 & 0.973 & 0.493 & 0.300 & 0.269 & 0.004 & 0.000 \\
\bottomrule
\end{tabular}
\end{table}

\begin{table}[t]\centering
\caption{Wang CascadedGCNN \citep{wang2025bridging}, retrained under their code
and recipe, weak versus strong attacks. Their protocol (normalized $[-1,1]$
space, their eps $0.03$ $=$ pixel $0.015$) is reported in the top block; the
pixel-space threat-matched AutoAttack (deterministic APGD-CE$+$APGD-T) in the
bottom. FGSM at their headline budget reproduces their numbers; AutoAttack erases
the robustness. Source:
\texttt{paper/results/b1\_wang/cascaded\_s0\_eval.json}
(their\_protocol / pixel\_space blocks, $n=1000$).}
\label{tab:wang-weak}
\footnotesize
\begin{tabular}{lccccc}
\toprule
\multicolumn{6}{l}{\emph{Their protocol} (normalized eps; clean $=87.2$)}\\
budget (norm.) & $0.01$ & $0.02$ & $0.03$ & $0.04$ & $0.05$ \\
FGSM   & 42.27 & 34.47 & 28.53 & 24.08 & 21.41 \\
PGD-40 & 20.24 &  5.53 &  1.75 &  0.74 &  0.34 \\
\midrule
\multicolumn{6}{l}{\emph{Pixel-space threat-matched} (clean $=88.6$)}\\
budget (pixel) & $0.01$ & $0.015$ & $0.03$ & & \\
PGD-40 (ref)   & 6.4 & 1.0 & 0.2 & & \\
AutoAttack (APGD) & 2.7 & 0.1 & 0.0 & & \\
\bottomrule
\end{tabular}
\end{table}

\subsection{Per-arm adversarial-training robust accuracies}\label{app:at-tables-full}
Under adversarial training the exactly invariant (circular / exact-cyclic) arm,
whose shift-consistency is $1.0$ by construction, is the least robust in every
width, while the threat-matched $\etaL$ orders robustness. Table~\ref{tab:at-cifar-full}
gives the CIFAR-10 $\ell_\infty$ and $\ell_2$ per-arm tables and
Table~\ref{tab:at-mf-full} the MNIST and Fashion-MNIST $\ell_\infty$ tables. The
four arms share a channel schedule and so an identical parameter count at each
width.

\begin{table}[t]\centering
\caption{Capacity-matched CIFAR-10 dissection under PGD adversarial training,
per-arm means. $\ell_\infty$ block: PGD-AT at $\varepsilon{=}8/255$ ($7$ steps),
evaluated by AutoAttack at $8/255$; $\eta/L_1$ is the threat-matched ($\ell_1$
dual) ratio. $\ell_2$ block: PGD-AT at $\varepsilon{=}0.5$ ($10$ steps),
AutoAttack at $\ell_2{=}0.5$; $\eta/L_2$ is the matched ratio. Sources:
\texttt{paper/results/at\_cifar\_linf\_full10k\_20260623\_185505.json} and
\texttt{at\_cifar\_l2\_full10k\_20260624\_111942.json} (full $10$k test set,
$2$ seeds).}
\label{tab:at-cifar-full}
\footnotesize
\begin{tabular}{ll ccc cc c}
\toprule
arm & $w$ & clean & consist. & robust acc. & $\eta$ & $L_1$ & $\eta/L_1$ \\
\midrule
\multicolumn{8}{l}{\emph{CIFAR-10}, $\ell_\infty{=}8/255$ (AutoAttack)}\\
standard     & 32 & 0.712 & 0.794 & 0.375 & 1.48 & 35.3 & 1.190 \\
anti-aliased & 32 & 0.724 & 0.802 & 0.388 & 1.56 & 36.3 & 1.191 \\
exact cyclic & 32 & 0.569 & 1.000 & 0.261 & 0.80 & 21.6 & 1.104 \\
shift-aug    & 32 & 0.694 & 0.867 & 0.359 & 1.35 & 32.3 & 1.150 \\
standard     & 64 & 0.766 & 0.820 & 0.408 & 1.89 & 46.1 & 1.133 \\
anti-aliased & 64 & 0.775 & 0.823 & 0.419 & 1.96 & 46.4 & 1.154 \\
exact cyclic & 64 & 0.633 & 1.000 & 0.294 & 1.05 & 28.8 & 1.081 \\
shift-aug    & 64 & 0.750 & 0.891 & 0.395 & 1.72 & 42.1 & 1.100 \\
\midrule
arm & $w$ & clean & consist. & robust acc. & $\eta$ & $L_2$ & $\eta/L_2$ \\
\midrule
\multicolumn{8}{l}{\emph{CIFAR-10}, $\ell_2{=}0.5$ (AutoAttack)}\\
standard     & 32 & 0.804 & 0.829 & 0.581 & 2.40 & 2.45 & 0.981 \\
anti-aliased & 32 & 0.815 & 0.832 & 0.594 & 2.52 & 2.48 & 1.016 \\
exact cyclic & 32 & 0.711 & 1.000 & 0.464 & 1.54 & 1.79 & 0.865 \\
shift-aug    & 32 & 0.794 & 0.902 & 0.570 & 2.24 & 2.28 & 0.984 \\
standard     & 64 & 0.846 & 0.847 & 0.609 & 3.16 & 3.28 & 0.964 \\
anti-aliased & 64 & 0.854 & 0.857 & 0.628 & 3.23 & 3.25 & 0.995 \\
exact cyclic & 64 & 0.769 & 1.000 & 0.511 & 1.95 & 2.31 & 0.843 \\
shift-aug    & 64 & 0.841 & 0.924 & 0.613 & 2.86 & 2.88 & 0.993 \\
\bottomrule
\end{tabular}
\end{table}

\begin{table}[t]\centering
\caption{MNIST and Fashion-MNIST dissection under PGD adversarial training
($\ell_\infty$), per-arm means. MNIST evaluated at $\ell_\infty{=}0.3$,
Fashion-MNIST at $\ell_\infty{=}0.1$, both by AutoAttack. The exactly invariant
arm has consistency $1.0$ and the lowest robust accuracy in each width. Sources:
\texttt{paper/results/at\_mnist\_linf\_full10k\_20260625\_043653.json} and
\texttt{at\_fashion\_linf\_full10k\_20260625\_160748.json}.}
\label{tab:at-mf-full}
\footnotesize
\begin{tabular}{ll ccc}
\toprule
arm & $w$ & clean & consist. & robust acc. \\
\midrule
\multicolumn{5}{l}{\emph{MNIST}, $\ell_\infty{=}0.3$}\\
standard     & 32 & 0.994 & 0.945 & 0.942 \\
anti-aliased & 32 & 0.994 & 0.953 & 0.942 \\
exact cyclic & 32 & 0.994 & 1.000 & 0.926 \\
shift-aug    & 32 & 0.994 & 0.998 & 0.939 \\
standard     & 64 & 0.994 & 0.946 & 0.947 \\
anti-aliased & 64 & 0.994 & 0.954 & 0.947 \\
exact cyclic & 64 & 0.994 & 1.000 & 0.934 \\
shift-aug    & 64 & 0.995 & 0.999 & 0.952 \\
\midrule
\multicolumn{5}{l}{\emph{Fashion-MNIST}, $\ell_\infty{=}0.1$}\\
standard     & 32 & 0.876 & 0.684 & 0.776 \\
anti-aliased & 32 & 0.876 & 0.683 & 0.776 \\
exact cyclic & 32 & 0.855 & 1.000 & 0.724 \\
shift-aug    & 32 & 0.862 & 0.958 & 0.747 \\
standard     & 64 & 0.886 & 0.680 & 0.781 \\
anti-aliased & 64 & 0.885 & 0.690 & 0.783 \\
exact cyclic & 64 & 0.868 & 1.000 & 0.736 \\
shift-aug    & 64 & 0.870 & 0.964 & 0.760 \\
\bottomrule
\end{tabular}
\end{table}

\subsection{RobustBench anisotropy table}\label{app:robustbench-full}
Among already-robust models $\etaL$ no longer orders robustness; the gradient's
anisotropy $A=\lVert\grad M\rVert_1/\lVert\grad M\rVert_2$ does.
Table~\ref{tab:robustbench-full} gives the per-model panel for thirty official
RobustBench CIFAR-10 $\ell_\infty$ models: clean accuracy, robust accuracy
(APGD-CE $\varepsilon{=}8/255$, $100$ iters), the anisotropy $A$, its normalized
form $A/\sqrt\dd$, and $\eta/L_1$. The correlation summary follows the table.

\begin{table}[t]\centering
\caption{Thirty RobustBench CIFAR-10 $\ell_\infty$ models, sorted by robust
accuracy. $A$ is the mean gradient anisotropy, $\sqrt\dd{=}55.4$. Source:
\texttt{llm\_transfer/pilots/c1\_tower/robustbench\_val/results.json} and
\texttt{FINAL\_REPORT.txt} ($n{=}1000$ eval, $500$ correct points attacked). The
in-house APGD ranking agrees with official AutoAttack at Spearman $+0.96$.}
\label{tab:robustbench-full}
\scriptsize
\begin{tabular}{lrrrrr}
\toprule
model & clean & robust & $A$ & $A/\sqrt\dd$ & $\eta/L_1$ \\
\midrule
Standard                             & 0.948 & 0.000 & 35.54 & 0.641 & 0.00358 \\
Wong2020Fast                         & 0.842 & 0.451 & 25.18 & 0.454 & 0.04112 \\
Andriushchenko2020Understanding      & 0.805 & 0.485 & 26.43 & 0.477 & 0.04299 \\
Ding2020MMA                          & 0.847 & 0.501 & 24.00 & 0.433 & 0.02538 \\
Cui2020Learnable\_34\_10             & 0.885 & 0.522 & 21.68 & 0.391 & 0.07622 \\
Engstrom2019Robustness               & 0.887 & 0.527 & 21.34 & 0.385 & 0.04883 \\
Cui2020Learnable\_34\_20             & 0.894 & 0.549 & 21.23 & 0.383 & 0.06939 \\
Addepalli2021Towards\_RN18           & 0.814 & 0.554 & 21.82 & 0.394 & 0.05044 \\
Addepalli2022Efficient\_RN18         & 0.866 & 0.556 & 24.43 & 0.441 & 0.05918 \\
Huang2020Self                        & 0.833 & 0.561 & 20.37 & 0.367 & 0.07275 \\
Hendrycks2019Using                   & 0.887 & 0.562 & 20.97 & 0.378 & 0.05281 \\
Zhang2020Attacks                     & 0.846 & 0.570 & 21.31 & 0.384 & 0.06717 \\
Sehwag2021Proxy\_R18                 & 0.852 & 0.579 & 21.75 & 0.392 & 0.06051 \\
Sehwag2020Hydra                      & 0.901 & 0.586 & 21.47 & 0.387 & 0.07087 \\
Chen2021LTD\_WRN34\_10               & 0.863 & 0.592 & 20.29 & 0.366 & 0.06534 \\
Gowal2021Improving\_R18\_ddpm\_100m  & 0.881 & 0.604 & 21.36 & 0.385 & 0.05745 \\
Addepalli2022Efficient\_WRN\_34\_10  & 0.885 & 0.609 & 21.96 & 0.396 & 0.07148 \\
Jia2022LAS-AT\_70\_16                & 0.852 & 0.612 & 19.50 & 0.352 & 0.06579 \\
Sridhar2021Robust                    & 0.901 & 0.613 & 20.05 & 0.362 & 0.07828 \\
Sridhar2021Robust\_34\_15            & 0.874 & 0.614 & 18.30 & 0.330 & 0.08935 \\
Gowal2020Uncovering\_28\_10\_extra   & 0.899 & 0.624 & 20.09 & 0.362 & 0.08079 \\
Rade2021Helper\_ddpm                 & 0.882 & 0.624 & 20.57 & 0.371 & 0.06741 \\
Wu2020Adversarial\_extra             & 0.885 & 0.627 & 19.99 & 0.361 & 0.07011 \\
Huang2021Exploring\_ema              & 0.914 & 0.632 & 19.08 & 0.344 & 0.09382 \\
Rebuffi2021Fixing\_28\_10\_cutmix\_ddpm & 0.881 & 0.638 & 21.54 & 0.389 & 0.07073 \\
Pang2022Robustness\_WRN28\_10        & 0.886 & 0.640 & 20.61 & 0.372 & 0.06470 \\
Zhang2020Geometry                    & 0.898 & 0.659 & 19.88 & 0.359 & 0.04861 \\
Pang2022Robustness\_WRN70\_16        & 0.891 & 0.663 & 18.85 & 0.340 & 0.07432 \\
Rebuffi2021Fixing\_70\_16\_cutmix\_ddpm & 0.884 & 0.667 & 18.68 & 0.337 & 0.08127 \\
Rebuffi2021Fixing\_106\_16\_cutmix\_ddpm & 0.886 & 0.668 & 18.08 & 0.326 & 0.08224 \\
\bottomrule
\end{tabular}
\end{table}

\paragraph{Correlation summary (RobustBench, $n{=}30$).} Lower $A$ means more
robust: $\mathrm{Spearman}(A,\text{robust})=-0.79$ ($p=2.0\times10^{-7}$),
bootstrap $95\%$ interval $[-0.92,-0.56]$ ($5000$ resamples). It survives
controlling for clean accuracy (partial $-0.77$) and for $\eta/L_1$ itself
(partial $-0.65$). Restricted to the twenty-nine robust models (dropping the
single non-robust \emph{Standard}) the ranking is $-0.77$ (the jackknife drop-of-%
\emph{Standard} value). The size $\eta/L_1$ detects robustness but ranks it only
weakly ($\mathrm{Spearman}(\eta/L_1,\text{robust})=+0.61$). The cross-check
against official AutoAttack gives $\mathrm{Spearman}(A,\text{officialAA})=-0.79$,
partial controlling clean accuracy $-0.76$.

\subsection{Crown-jewel hardened seventeen-tower four-cell table}\label{app:crown-full}
On zero-shot CLIP over ImageNet-100 the two-by-two crosses modality (image, text)
with perturbation type (average-case invariance, worst-case adversarial). Both
invariance cells saturate; only the image-adversarial cell $S$ (full targeted
AutoAttack) carries cross-encoder spread; the threat-matched $\eta/L_1$ predicts
it while the consistencies do not. Table~\ref{tab:crown-full} gives the per-tower
panel and Table~\ref{tab:crown-corr} the pre-registered correlation table with
confidence intervals.

\begin{table}[t]\centering
\caption{Zero-shot CLIP, ImageNet-100, seventeen encoders, full targeted
AutoAttack at $\varepsilon{=}4/255$ ($0.01569$). ``rob.'' marks image-adversarial
robustness (robust$=$T for the thirteen adversarially robust towers).
SC${=}$shift-consistency, PC${=}$paraphrase-consistency, $S{=}$image-adversarial
accuracy (AutoAttack), $S_{\text{text}}{=}$worst-of-eleven-paraphrase accuracy.
The text-hardened tower \emph{leaf\_L\_text} (last row, marked~$\dagger$) has the
highest paraphrase-consistency tier despite a middling vision $\eta/L_1$; it is
counted among the four non-image-robust towers. Source:
\texttt{llm\_transfer/results/c1\_tower/crownjewel\_hardened.json}
($n_{\text{eval}}{=}2000$, $n_{\text{attack}}{=}400$).}
\label{tab:crown-full}
\scriptsize
\begin{tabular}{llccccc r}
\toprule
tower & rob. & clean & SC & PC & $S$ & $S_{\text{text}}$ & $\eta/L_1$ \\
\midrule
clip\_b32\_laion2b & F & 0.843 & 0.9598 & 0.9789 & 0.000 & 0.9015 & 0.00069 \\
fare4\_b32        & T & 0.691 & 0.9628 & 0.9690 & 0.393 & 0.8611 & 0.01756 \\
tecoa4\_b32       & T & 0.770 & 0.9658 & 0.9675 & 0.575 & 0.8500 & 0.02342 \\
fare4\_b16        & T & 0.760 & 0.9756 & 0.9623 & 0.495 & 0.8361 & 0.01963 \\
tecoa4\_b16       & T & 0.794 & 0.9770 & 0.9724 & 0.573 & 0.8741 & 0.02569 \\
clip\_b16\_openai & F & 0.844 & 0.9675 & 0.9808 & 0.000 & 0.8986 & 0.00049 \\
simclip4          & T & 0.836 & 0.9790 & 0.9738 & 0.540 & 0.8864 & 0.01245 \\
simclip2          & T & 0.877 & 0.9824 & 0.9794 & 0.340 & 0.9002 & 0.01124 \\
clip              & F & 0.900 & 0.9816 & 0.9837 & 0.000 & 0.9216 & 0.00031 \\
clip\_l14\_laion2b& F & 0.891 & 0.9855 & 0.9853 & 0.000 & 0.9321 & 0.00025 \\
fare4\_cnxt       & T & 0.803 & 0.9804 & 0.9726 & 0.507 & 0.8773 & 0.02151 \\
tecoa4\_cnxt      & T & 0.813 & 0.9737 & 0.9732 & 0.620 & 0.8937 & 0.02766 \\
fare4             & T & 0.853 & 0.9825 & 0.9756 & 0.580 & 0.8927 & 0.01220 \\
tecoa4            & T & 0.875 & 0.9895 & 0.9820 & 0.693 & 0.9143 & 0.01636 \\
fare2             & T & 0.885 & 0.9798 & 0.9818 & 0.320 & 0.9192 & 0.01130 \\
tecoa2            & T & 0.915 & 0.9901 & 0.9894 & 0.578 & 0.9448 & 0.02207 \\
leaf\_L\_text$^\dagger$ & F & 0.887 & 0.9811 & 0.9850 & 0.338 & 0.9234 & 0.01174 \\
\bottomrule
\end{tabular}
\end{table}

The four-cell saturation is explicit (means over the $17$ towers, standard
deviations reported as in the main text, sample convention): image-invariance
$\mathrm{SC}=0.977\pm0.009$ (range $[0.960,0.990]$), text-invariance
$\mathrm{PC}=0.977\pm0.007$ (range $[0.962,0.989]$); the image-adversarial cell
$S=0.385\pm0.244$ (range $[0,0.69]$, $27.6\times$ the SC spread by the analysis
file's population-sd ratio); the text-adversarial cell
$S_{\text{text}}=0.896\pm0.029$ (range $[0.836,0.945]$, $4.1\times$ the PC
spread).

\begin{table}[t]\centering
\caption{Pre-registered correlations on the seventeen-tower panel, with
$95\%$ intervals ($n{=}17$ unless noted). $\eta/L_1$ orders the image-adversarial
cell but not the text cell (raw anti-correlation is a clean-accuracy artifact,
vanishing under the partial); paraphrase-consistency co-moves with $S_{\text{text}}$
because both are average-versus-worst over the same narrow paraphrase family. The
per-image row pools the thirteen robust towers within-tower ($n{=}2600$). Source:
\texttt{llm\_transfer/results/c1\_tower/crownjewel\_hardened\_analysis.json}.}
\label{tab:crown-corr}
\small
\begin{tabular}{llc}
\toprule
relation & Spearman & $95\%$ CI \\
\midrule
$\eta/L_1$ vs $S$ (image-adversarial)          & $+0.80$ & $[+0.44,+0.95]$ \\
\quad partial, controlling clean accuracy       & $+0.82$ & $[+0.50,+0.95]$ \\
$\eta/L_1$ vs $S_{\text{text}}$ (raw)           & $-0.54$ & $[-0.87,-0.08]$ \\
\quad partial, controlling clean accuracy       & $+0.006$ & $[-0.55,+0.51]$ \\
PC vs $S_{\text{text}}$                         & $+0.98$ & $[+0.89,+1.00]$ \\
\quad PC vs $S_{\text{text}}$, partial | clean  & $+0.77$ & $[+0.20,+1.00]$ \\
SC vs $S$                                        & $+0.21$ & $[-0.37,+0.67]$ \\
PC vs $S$                                        & $-0.30$ & $[-0.75,+0.22]$ \\
per-image $\eta/L_1$ vs radius (pooled robust, $n{=}2600$) & $+0.89$ & $[+0.88,+0.90]$ \\
\bottomrule
\end{tabular}
\end{table}

Within the seventeen-tower panel, per-tower the per-image $\eta/L_1$ vs radius
Spearman ranges from $+0.52$ (clip\_b16\_openai, a floor-dominated non-robust
tower) up to $+0.97$ (fare4\_b32), with the pooled-robust value $+0.89$; this is
the hardened panel and is separate from the ten-robust-encoder set of
\S\ref{sec:etaL} (per-encoder $+0.74$ to $+0.96$). The single text-hardened tower
\emph{leaf\_L\_text} has
among the highest paraphrase-consistency and $S_{\text{text}}$ (PC $0.9850$,
$S_{\text{text}}$ $0.9234$) despite a middling vision $\eta/L_1$ ($0.01174$),
placing text robustness as a separate axis raised by text training and not by the
vision margin.

\subsection{MNIST and Fashion-MNIST anisotropy sweep tables}\label{app:anisosweep-full}
On low-dimensional families a fixed training budget compresses $A$ into a narrow
band where it does not rank; an adversarial-training-strength sweep widens it and
restores the ranking law. Table~\ref{tab:anisosweep-corr} gives the twenty-four
cell per-dataset correlations at each evaluation budget. Over the sweep $A$ spans
$[2.38,16.59]$ on MNIST and $[3.99,13.17]$ on Fashion, and is collinear with
training strength ($\mathrm{Spearman}(A,\varepsilon_{\mathrm{train}})\approx-0.85$
MNIST, $-0.81$ Fashion), so the diverse RobustBench panel remains the cleaner
cross-model evidence.

\begin{table}[t]\centering
\caption{Anisotropy $A$ versus threat-matched $\ell_\infty$ robust accuracy over
$24$ cells per dataset (six training budgets $\times$ two widths $\times$ two
seeds), with the partial controlling clean accuracy and, for contrast, the
mismatched $\ell_2$ (DDN) radius. $A$ ranks robustness in the training norm at the
informative budgets; the sign is correct only in the training norm (on the
mismatched $\ell_2$ radius it is right on Fashion but wrong on MNIST, collapsing
to null under the clean-accuracy control). Sources:
\texttt{paper/results/anisosweep\_analysis\_mnist\_linf.json} and
\texttt{...\_fashion\_linf.json} (per-cell rows in
\texttt{at\_\{mnist,fashion\}\_linf\_anisosweep\_*.json}).}
\label{tab:anisosweep-corr}
\small
\begin{tabular}{llcc}
\toprule
dataset / metric & budget & $\mathrm{Spearman}(A,\cdot)$ & partial | clean \\
\midrule
\multicolumn{4}{l}{\emph{MNIST} ($A\in[2.4,16.6]$)}\\
AutoAttack $\ell_\infty$ & $\varepsilon{=}0.2$ & $-0.83$ & $-0.75$ \\
AutoAttack $\ell_\infty$ & $\varepsilon{=}0.3$ & $-0.89$ & $-0.74$ \\
$\ell_2$ radius (DDN, mismatched) & n/a & $+0.48$ & $+0.09$ \\
\midrule
\multicolumn{4}{l}{\emph{Fashion-MNIST} ($A\in[4.0,13.2]$)}\\
AutoAttack $\ell_\infty$ & $\varepsilon{=}0.1$  & $-0.61$ & $-0.53$ \\
AutoAttack $\ell_\infty$ & $\varepsilon{=}0.15$ & $-0.80$ & $-0.67$ \\
$\ell_2$ radius (DDN, mismatched) & n/a & $-0.74$ & $-0.26$ \\
\bottomrule
\end{tabular}
\end{table}

The mismatched-$\ell_2$ sign nuance is that strong high-$\varepsilon$
$\ell_\infty$ training over-fits the $\ell_\infty$ ball and shrinks the $\ell_2$
radius, so on MNIST the raw $A$ vs $\ell_2$ radius correlation is $+0.48$
(``wrong'' sign relative to the $\ell_\infty$ story) and collapses to $+0.09$
once clean accuracy is controlled; on Fashion the mismatched radius keeps the
correct sign ($-0.74$ raw). Every correlation was recomputed three ways with no
gradient masking; the effective-dimension deficit
$r\approx(\eta/L_1)/(1+CA)$ this evidence supports remains, as the main text
scopes it, a proposed and well-controlled mechanism rather than a network-level
theorem.

\subsection{The diffusion data-axis, TM-regression interventions, and the
certificate-count mechanism}\label{app:extra-full}
\paragraph{The diffusion data-axis.} Intervening on the training \emph{data}
instead of the architecture, through diffusion-augmented adversarial training,
obeys the same threat-matched law, an independent axis for the diagnostic. Across
adversarial-training budgets $\varepsilon\in\{2,4,8,16\}/255$ the raw
threat-matched ratio anti-correlates with robust accuracy ($-0.95$) while the
budget-normalized ratio $\eta/(L\varepsilon)$, the certified radius in units of
the attack, collapses the sweep onto one monotone curve (Pearson $+0.94$); within
a fixed budget, which is how the dissections compare models, the raw ratio is the
right form. At matched $\eta/L_1$ the diffusion-data models sit measurably above
the real-only models ($+0.10$ AutoAttack, $t\approx+5$ in the budget sweep), an
offset that tracks the train--test robust gap the data removes and vanishes at
$\varepsilon=16/255$ where there is no robust overfitting left to prevent. The
mean test-set ratio is therefore not a sufficient statistic across interventions;
the per-sample certified fraction is (Proposition~\ref{prop:count-app}), which
localizes the limitation to the mean summary rather than to the certificate.

\paragraph{Four interventions confirm $\etaL$ is not a trainable target.} The
scale-invariance of Lemma~\ref{lem:ratiodegen} makes the ratio objective
degenerate and the radial floor of Proposition~\ref{prop:split-app} ties the
gradient to the margin; four controlled interventions on a standard backbone
settle the empirical question. \emph{(1)~Isotropic gradient penalty} on the
threat-matched margin gradient $\lVert\grad M\rVert_1$: only the threat-matched
$\ell_1$ penalty raises $\etaL$ (the input-gradient $\ell_1$ norm falls from $38$
to $22$), while the mismatched $\ell_2$ and loss-gradient penalties
\citep{ross2018improving,finlay2021scaleable} leave it unchanged; but AutoAttack
robust accuracy is statistically unchanged for every moderate penalty (paired
$p>0.4$) and falls under the strongest matched penalty ($-0.011$, $p=0.026$), with
no gradient masking, because penalizing the gradient lowers the margin alongside
it ($\eta$ falls from $1.58$ to $1.03$). \emph{(2)~Ratio maximization} of
$\lVert\grad M\rVert_1/M$ drives the classifier to the trivial constant solution
(the scale-invariance degeneracy). \emph{(3)~Anisotropic hinge penalty} and
\emph{(4)~directional penalty} at the adversarial point collapse in the same way:
clean accuracy falls from $0.728$ to $0.38$ and $0.29$, AutoAttack from $0.390$
to $0.258$ and $0.201$, and the margin from $1.59$ to $0.24$ and $0.05$, every
quantity dropping together with no gradient masking. All four collapse the margin.

\paragraph{The certificate-count mechanism.} In the locally affine regime the
first-order radius $r_1(x)=M(x)/\lVert\grad M(x)\rVert_1$ is an \emph{exact}
$\ell_\infty$ survival test, so robust accuracy is a lower-tail functional of the
$r_1$ distribution.

\begin{proposition}[Exact first-order certificate count on affine cells]\label{prop:count-app}
Fix $\varepsilon>0$ and suppose $M$ is affine on $B_\infty(x,\varepsilon)$, so
$M(x+\delta)=M(x)+\langle g,\delta\rangle$ there with $g=\grad M(x)$. Then
$\inf_{\lVert\delta\rVert_\infty\le\varepsilon}M(x+\delta)=M(x)-\varepsilon\lVert g\rVert_1$,
so $x$ survives every $\ell_\infty$ perturbation of radius $\varepsilon$ iff
$M(x)\ge\varepsilon\lVert g\rVert_1$; equivalently (with $r_1=\infty$ when $g=0$
and $M>0$) iff $M(x)\ge0$ and $r_1(x)\ge\varepsilon$. If the model is affine on
$B_\infty(x,\varepsilon)$ for almost every test point, the population robust
accuracy equals the certified fraction,
$\mathrm{RobAcc}(\varepsilon)=\mathbb P_x(M(x)\ge0,\ r_1(x)\ge\varepsilon)$.
\end{proposition}
\begin{proof}
By affineness,
$\inf_\delta M(x+\delta)=M(x)+\inf_{\lVert\delta\rVert_\infty\le\varepsilon}\langle g,\delta\rangle$.
H\"older gives $\langle g,\delta\rangle\ge-\varepsilon\lVert g\rVert_1$, attained
by the feasible $\delta^\star_j=-\varepsilon\,\sign(g_j)$; survival on the whole
ball is nonnegativity of the infimum, i.e.\ $M(x)\ge\varepsilon\lVert g\rVert_1$,
which for $g\neq0$ reads $r_1(x)\ge\varepsilon$ and for $g=0$ reduces to
$M(x)\ge0$. Taking probability over the test distribution gives the count.
\end{proof}

\begin{remark}[Reading]\label{rem:count-app}
The forward implication is the margin-over-local-Lipschitz certificate
specialized to the $\ell_\infty/\ell_1$ dual pair
\citep{hein2017formal,tsuzuku2018lipschitz}; the content is the converse: on an
affine cell the perturbation $-\varepsilon\,\sign(\grad M)$ attains the dual norm
exactly, so the certificate is an exact survival test, not merely a lower bound.
For a multi-class score $M$ is a minimum of pairwise affine margins, so ``$M$
affine on the ball'' says the active rival does not switch there; the $+0.998$
agreement between the certified fraction and AutoAttack indicates this regime
holds at the budgets used, and its breakdown at larger budgets or curvier models
is the open edge. The budget-normalized mean law is the scale-family shadow of
this count, and the data-axis offset is a tail effect: the certified fraction
absorbs the dose coefficient ($+0.053\to+0.006$).
\end{remark}
